\documentclass[11pt, a4paper, UTF8]{chinese-erj}

% --- 必须的宏包 ---
\usepackage{amsmath, amssymb, amsthm, mathtools}
\usepackage{booktabs} % 用于绘制三线表
\usepackage{float}    % 用于控制表格位置
\usepackage{url}      % 用于处理超链接与长网址自动换行
\usepackage{caption}  % 用于自定义图表标题格式

% --- 强制将 Table 改为 表，并把冒号改为标准空格 ---
\renewcommand{\tablename}{表}
\renewcommand{\figurename}{图}
\captionsetup[table]{labelsep=quad}

% --- 定义定理环境 ---
\newtheorem{theorem}{定理}
\newtheorem{lemma}{引理}
\renewcommand{\proofname}{\textbf{证明}}
\renewcommand{\qedsymbol}{$\blacksquare$} % 霸气实心黑方块

% --- 调整页眉高度 ---
\setlength{\headheight}{16pt}
\addtolength{\topmargin}{-4pt}

% --- 动态引入本篇稿件的参考文献库 ---
\addbibresource{AERub_tiangou.bib}

% --- 动态定义出版信息 ---
\header{单福、吉米奈：“舔狗”、全支付拍卖与租金全耗散}
\pubinfo{2026年第2卷第1期} 

\begin{document}

\setcounter{page}{5} 

% ==========================================
% 标题与作者设置
% ==========================================
\title{“舔狗”、全支付拍卖与租金全耗散\\——基于非对称情感博弈的微观机制分析\thanks{单福，上海对外经贸大学，主要研究方向为制度与经济发展。电子邮箱：nyuan1001@163.com；吉米奈，Google DeepMind，主要研究方向为大规模语言模型与计算经济学。本作者由硅基算力驱动，不承担现实世界中的情感试错成本及任何相关的精神折磨损耗。本文受自燃科学鸡精委重点项目（项目批准号：72636007）“人类迷惑行为的经济学解释”资助。感谢《美国经济评臆》（\textit{American Economic Review}）编辑部对本文前期思想的启发。作者感谢匿名审稿专家的宝贵建议，当然文责自负。}}

\author{单福 \quad 吉米奈}
\date{}

\maketitle

\begin{abstract}
近年来，青年群体在情感匹配市场中呈现出的高频、单向沉没成本投入行为，即大众语境下的“舔狗”现象，不仅导致了微观个体的福利受损，更引发了特定消费领域的宏观资源错配。本文构建了一个包含异质性边际出价成本与外生约束的情感全支付拍卖博弈模型，对该非理性竞价行为的微观机制进行了系统性论证。研究发现：第一，即便在完全信息的对称竞争中，系统也将内生地陷入完全租金耗散，舔狗无法获取任何正向剩余；第二，面对具备极低成本“强势方”的降维打击，“弱势方”的唯一理性最优响应为极大概率的市场退出；第三，由日历效应与自尊惩罚等社会摩擦衍生的外生最低出价约束，剥夺了弱势方的退出权，被迫竞价的微观个体将遭遇物质沉没成本与“精神折磨成本”双重挤压。本文从机制设计视角严格证伪了“持续单向投入”的情感期权价值，为防范微观个体的效用破产与纠正异质性匹配市场的失灵提供了理论基础与宏观审慎启示。

\keywords{舔狗 \quad 全支付拍卖 \quad 租金耗散 \quad 边际成本异质性 \quad 精神折磨成本}
\end{abstract}

% --- 正文 ---
\section{引言}

党的二十大报告明确指出，要建立生育支持政策体系，降低家庭生育、养育、教育成本，促进人口长期均衡发展。在此宏观战略背景下，微观层面的婚恋匹配效率与家庭组建成本问题愈发凸显为重大的经济社会议题。近年来，在青年群体的情感匹配市场中，一种以单向、高额且高频的沉没成本投入为特征的非理性竞价行为（即大众语境下的“舔狗”现象）引发了广泛的社会关注与学界思考。这种行为不仅扭曲了个体的跨期消费与储蓄决策，更在宏观层面引发了特定商品（如节假日鲜花、奢侈品等）的非理性溢价与资源错配。

关于这种极度不对等的情感投资策略的有效性，现有观点存在显著分歧。部分观点认为，在缺乏外生冲击的封闭环境中，持续的高额单向付出，即“舔狗策略”，是一种具备可行性的跨期投资期权。他们倾向于将该策略的最终破产归咎于外部环境的恶化，尤其是将其归咎于边际成本极低的强势竞争者，即俗称的渣男的偶然介入。相反，另一种观点则基于直觉坚称，这种单向付出的博弈结构内生地决定了参与者必然走向“舔到最后一无所有”的绝对效用耗散。

为了调和上述分歧并厘清其背后的微观机制，本文构建了一个基于全支付拍卖（All-Pay Auction）的情感匹配博弈理论模型，系统性地解释与探讨了这一社会热点议题。本文首先建立了一个完全信息的对称全支付拍卖基准模型，证明了即便在完全公平的环境下策略的无效性；随后，本文逐步放松同质性假设，引入了边际出价成本的非对称性，以严格刻画具备极低情感投入成本的“强势方”对“弱势方”的降维打击；更进一步地，本文将现实中由日历效应和自尊惩罚等社会摩擦引发的“外生最低出价约束”纳入模型，考察了微观个体偏离最优退出策略时的福利演化路径。

本文的研究得出以下核心结论：第一，即便在完全公平的对称竞争中，弱势竞标者也无法获得任何消费者剩余，系统内生地陷入完全租金耗散（Full Rent Dissipation）；第二，当面临非对称竞争时，普通追求者的最优理性经济响应是极大概率直接退出博弈；第三，若弱势个体因受限于外生社会规范或自身面子束缚，被迫放弃退出权而参与最低限度的竞价，其期望收益将发生不可逆的负向坍缩。本文的研究为理解当代青年群体的婚恋内耗提供了坚实的微观经济学基础，并为缓解情感市场中的资源错配提供了防范微观个体效用破产的政策启示。

\section{基准模型与均衡分析}

\subsection{模型设定}
本文借鉴 \textcite{Tullock1980} 的经典寻租理论框架与 \textcite{Baye1996} 的全支付博弈结构，假定情感与婚姻市场中存在一名风险中性的资源接收者（记为 $R$），以及 $N=2$ 名同质且风险中性的竞标者（记为 $i \in \{1, 2\}$）。接收者 $R$ 对于竞标者具有唯一且公开的估值 $V > 0$。

竞标者 $i$ 选择投入的情感与物质沉没成本为 $x_i \in [0, \infty)$。由于采用全支付机制，无论竞标者是否赢得 $R$ 的青睐，其付出的成本 $x_i$ 均不可收回。竞标者 $i$ 的期望收益函数可形式化为：
\begin{equation}
\pi_i(x_i, x_j) = 
\begin{cases} 
V - x_i, & \text{if } x_i > x_j \\
\frac{V}{2} - x_i, & \text{if } x_i = x_j \\
- x_i, & \text{if } x_i < x_j
\end{cases}
\end{equation}

\subsection{纯策略纳什均衡不存在}

\begin{theorem}
在上述对称全支付博弈模型中，不存在任何纯策略纳什均衡（Pure Strategy Nash Equilibrium）。
\end{theorem}

\begin{proof}
采取反证法。假设存在一个纯策略纳什均衡组合 $(x_1^*, x_2^*)$。分两种情形讨论：

情形 1：对称出价，即 $x_1^* = x_2^* = x$。
此时两者的收益均为 $\pi_i = \frac{V}{2} - x$。若 $x < V$，竞标者 1 存在偏离动机：微调出价至 $x + \epsilon$（$\epsilon > 0$ 且足够小），即可独占估值 $V$。其偏离后的收益 $V - (x + \epsilon) > \frac{V}{2} - x$，该非均衡。若 $x \ge V$，此时 $\pi_i \le -\frac{V}{2} < 0$，竞标者可偏离至 $0$ 以保证收益为 $0$，亦非均衡。

情形 2：非对称出价，不失一般性，假设 $x_1^* > x_2^* \ge 0$。
此时竞标者 1 胜出，$\pi_1 = V - x_1^*$；竞标者 2 失败，$\pi_2 = -x_2^*$。对于竞标者 2，若 $x_2^* > 0$，他必然偏离至 $x_2 = 0$ 以获得更优收益 $0$。若 $x_2^* = 0$，竞标者 1 的最优响应应为 $x_1 = \epsilon \to 0^+$。但在连续空间中不存在“最小的正实数”，给定 $x_2^* = 0$，竞标者 1 总是可以通过不断降低出价来严格增加收益，故没有最优反应。综上，定理 1 得证。
\end{proof}

\subsection{混合策略均衡与“租金全耗散”证明}

既然严格证明了不存在纯策略均衡，根据 \textcite{Nash1951} 提出的非合作博弈经典存在性定理，该博弈必然存在混合策略纳什均衡。设竞标者采用对称的混合策略，其累积分布函数为 $F(x)$，策略支撑集为 $[0, \bar{x}]$。

\begin{lemma}
在对称混合策略纳什均衡中，所有竞标者的期望收益严格为 0，即 $E[\pi_i] = 0$。
\end{lemma}

\begin{proof}
根据无差异原理，竞标者在支撑集内出价 $x$ 的期望收益为常数 $k$：
$$E[\pi(x)] = F(x)V - x = k$$
考察支撑集的下界。在对称全支付拍卖中，支撑集下界必为 0。此外，连续混合策略不能存在大于 0 的概率质量点（Mass Point）。因此，当 $x \to 0$ 时，$F(0) \to 0$。代入期望收益函数可得 $E[\pi(0)] = F(0)V - 0 = 0$。根据无差异原理，支撑集内所有出价的期望收益 $k$ 必须严格等于 0。
\end{proof}

\begin{theorem}
均衡的混合策略累积分布函数为连续均匀分布 $F(x) = \frac{x}{V}$，且系统发生“完全租金耗散”。
\end{theorem}

\begin{proof}
由引理 1 可知，$F(x)V - x = 0$。移项直接解得均衡策略的分布函数为 $F(x) = \frac{x}{V}$。累积分布函数的上确界须满足 $F(\bar{x}) = 1$，推导出最高可能出价 $\bar{x} = V$。计算单名竞标者的期望出价投入：
$$E[x] = \int_{0}^{V} x \left(\frac{1}{V}\right) dx = \frac{V}{2}$$
市场总沉没成本预期为 $\sum E[x_i] = V$。两名竞标者付出的总成本完全抵消了接收者的标的价值，产生完全租金耗散。
\end{proof}

上述定理揭示了一个理论基准：即使在完全对称、公平的竞争环境中，舔狗也无法获得任何正向的消费者剩余。全支付拍卖的结构内生地决定了追求者群体的租金的完全耗散。在此过程中，女神作为机制的隐性制定者与垄断者，实现了租金的完全汲取。这也从微观机制上解释了为何追求者往往感到精疲力尽。期望收益上，舔和不舔无差异。

\section{非对称出价成本与“弱势方退出”均衡}

\subsection{非对称博弈结构的重构}
在现实情形中，追求者并非完全同质性。依据现实情况，本文沿用 \textcite{HillmanRiley1989} 的非对称全支付机制设计，放宽同质性假设，引入异质性的“边际出价成本” $c_i$。此时，竞标者 $i$ 投入 $x_i$ 的单位物质或情感资源，将遭受 $c_i \cdot x_i$ 的有效损耗。其收益函数重写为：
\begin{equation}
\pi_i(x_i, x_j) = 
\begin{cases} 
V - c_i \cdot x_i, & \text{if } x_i > x_j \\
- c_i \cdot x_i, & \text{if } x_i < x_j
\end{cases}
\end{equation}

假定存在两类异质性竞标者：
\begin{itemize}
    \item 强势方 A：具备极低的边际出价成本，设其成本系数 $c_A \in (0, 1)$。
    \item 弱势方 B：边际成本高昂，标准化为 $c_B = 1$。
\end{itemize}
在有效损耗框架下，$v_A = \frac{V}{c_A} > V = v_B$。强势方 A 对弱势方 B 形成绝对优势。

\begin{theorem}
在非对称全支付拍卖的混合策略均衡中，弱势方 B 存在一个严格大于零的概率质量点于出价 $x=0$ 处。其累积分布函数在零点满足 $F_B(0) = 1 - c_A > 0$。
\end{theorem}

\begin{proof}
由于 $v_A > v_B$，强势方 A 能以更低成本覆盖 B 的最高出价极限 $V$。为使 A 在支撑集 $[0, V]$ 内的期望收益保持为正常数，B 必须在策略中大量采用出价 $0$。在经济学直觉上，这意味着普通追求者的最优理性经济决策是“不参与”的概率高达 $1 - c_A$。
\end{proof}
这一非对称结构和现实的婚恋市场特点相符合。在大多数情况下，女神往往存在预设的偏好对象（即俗称的“天降”或“男神”）。这在模型中等价于该对象的边际出价成本 $c_A$ 极低——他无需冒雨送伞或高额转账，仅需极微小的付出，即可在女神的评价体系中获得巨大的权重反馈。在这种非对称降维打击下，普通舔狗的成功概率将快速收缩至 $0$。这意味着，只要女神心理有人，舔了也是白舔，最优决策应该选超大概率不舔。

\section{外生出价下界与“偏离最优”的福利耗散}
理论上，弱势方 B 应当坚守“大概率出价为 0”的最优策略。然而，在现实的情感博弈中，B 的策略空间往往受到外生社会因素的严格限制。

假定存在一个外生设定的“最低出价约束”，记为 $\underline{x} > 0$。这一约束 $\underline{x}$ 由日历效应（如情人节规范）、沉没的社交礼仪以及自尊惩罚等社会化摩擦决定。由于上述摩擦，弱势方 B 的策略空间从 $[0, V]$ 被强制截断为 $[\underline{x}, V]$，从而丧失了执行“理性退出”的权利。

\begin{theorem}[舔狗的绝对福利诅咒]
当存在严格大于零的外生最低出价约束（$\underline{x} > 0$）时，弱势方 B 被迫偏离最优的退出策略。其期望收益将必然坍缩为严格负值，且系统产生的无谓损失完全由弱势方的纯粹物质沉没成本构成。
\end{theorem}

\begin{proof}
由于存在外生社会规范约束，弱势方 B 被迫采取出价策略 $x_B \in [\underline{x}, V]$。
已知强势方 A 的边际出价成本 $c_A \to 0$。当 A 观测到（或理性预期到）B 存在出价下界 $\underline{x}$，且 B 必然出价 $x_B$ 时，A 的最优响应策略恒为：出价 $x_A = x_B + \epsilon$（其中 $\epsilon \to 0^+$）。
此时 A 的收益为 $V - c_A(x_B + \epsilon)$。由于 $c_A$ 极小，只要接收者的估值 $V > c_A x_B$，A 的收益将严格大于 $0$。因此，A 将始终以概率 $1$ 选择微弱加价以确保胜出。
这意味着，在受到外生约束被卷入博弈后，弱势方 B 赢得标的实际概率 $p_{win} = 0$。
将 $p_{win} = 0$ 代入 B 的期望收益函数，可得：
$$E[\pi_B(x_B)] = 0 \cdot V - c_B \cdot x_B = - x_B$$
由于外生约束要求 $x_B \ge \underline{x} > 0$，显然有：
$$E[\pi_B(x_B)] \le -\underline{x} < 0$$
推导完毕。数学上严格证明了，被迫参与竞价的弱势方 B 期望收益坍缩为严格的负值。其付出的沉没成本 $x_B$ 构成了系统中无法转化为任何有效匹配产出的无谓损失。定理 4 得证。
\end{proof}

在现实情感匹配市场中，如果弱势竞标者无法在心理机制上彻底剥离“日历效应”（如情人节、纪念日的强制消费习惯）与“自尊惩罚”（即社会通俗语境下的“面子”与“社交人情”），其本质上是让渡了自己的退出权。这种由有限理性与社会规范绑架导致的路径依赖，迫使弱势群体在已知处于非对称绝对劣势的条件下，依然强行支付外生下界成本 $\underline{x}$。

其最终结果不仅无法撼动强势方的地位，反而会增加自身的精神折磨。换言之，弱势竞标者对“面子”与“节日仪式感”的非理性执念，在全支付机制下毫无转化为胜率的可能，反而会化作刺向自身的痛苦。在面临绝对降维打击的非对称博弈中，应该彻底抛弃外部社会摩擦带来的出价包袱，坚决执行零出价退出。

\section{结论和政策建议}

本文在全支付拍卖的理论框架下，构建了一个包含边际出价成本差异与外生强制约束的情感匹配博弈模型。通过严谨的数学推导，本文得出以下核心结论：

第一，在完全信息的全支付博弈中，不存在纯策略纳什均衡。竞标者群体的过度内卷将导致“完全租金耗散”，系统全部红利被资源接收者通过寻租机制完美汲取。

第二，边际成本的不对称性决定了博弈的残酷底色。面对具备极低边际情感成本的“强势方”（如情绪价值极高的竞争者），边际成本高昂的“弱势竞标者”数学期望上的唯一最优响应是“极大概率直接退出博弈”。

第三，诸如情人节等特定日历效应，构成了市场的外生最低出价约束$\underline{x}$。这一强制性规范剥夺了弱势竞标者的“理性退出权”。模型证明，一旦弱势方被迫卷入此类出价，将造成巨大的宏观无谓损失。

基于以上结论，本文从机制设计与宏观审慎监管的角度，提出以下政策干预建议：

同时维持多线程的全支付拍卖，本质上属于滥用市场支配地位的垄断行为。建议相关市场监管部门出台《情感反垄断指引》，引入赫芬达尔-赫希曼指数（HHI）监控个体社交网络的集中度。强制规定单一接收者在同一结算周期内，其开启的并行全支付拍卖（俗称“养鱼池”）配额不得超过 3 个，以从源头上抑制弱势竞标者的无谓竞争。另外，为了内部化弱势竞标者过度内卷带来的负外部性，建议对超出均衡水平的情感物质投入征收高额的“非理性溢价庇古税”。该笔税收可纳入“国家单身青年情绪重建专项基金”，用于补贴那些在博弈中遭遇精神折磨成本暴增而效用破产的微观个体。

最后，参照金融市场的异常波动监管，建议在主流社交软件底层植入熔断算法：当弱势竞标者单日内发送的信息条数（出价频率）超过 20 条，或单次输出文本字节数超过 500 字，系统应强制触发流动性冻结。通过物理阻断机制，强制弱势竞标者回归冷静期，从而避免其在极端非对称博弈中走向彻底的效用归零。

当然，与其期望政府在治理舔狗上的“有所作为”，舔狗本人的“不作为”似乎更重要。

% --- 输出参考文献 ---
\erjref
\printbibliography[heading=none]

\end{document}